% =====================================================================
% Section 4.X: Complexity analysis (formal)
%
% Replaces the brief "Complexity" paragraph in paper_main.tex with a
% dedicated subsection that:
%   - tabulates per-operation cost
%   - derives the closed-form total cost
%   - gives concrete operation counts for GEANT
%   - compares against the Dijkstra baselines
% =====================================================================

\subsection{Complexity}
\label{subsec:complexity}

We analyze the time and space complexity of
Algorithm~\ref{alg:core} per routing decision (one
source--destination pair). Let $n = |V|$, $e = |E|$, $\bar{d} =
2e/n$ the mean degree, $H = \lceil \alpha_{\text{budget}}
\cdot h_{\text{sp}}\rceil$ the hop budget (with $h_{\text{sp}}$
the shortest-path hop count from $s$ to $t$), and $B$ the number
of mass buckets. Recall that
$\alpha_{\text{budget}} = 1.25$ and $B = 32$ throughout this paper.

\paragraph{Per-operation cost.}
Table~\ref{tab:complexity} enumerates the four operations the DP
performs and the cost of each.

\begin{table}[h]
\centering
\caption{Per-operation cost of Algorithm~\ref{alg:core}.
The dominant term is the relaxation step, which is repeated
$H \cdot n \cdot \bar{d} \cdot B$ times.}
\label{tab:complexity}
\small
\begin{tabular}{lll}
\toprule
Operation & Frequency & Cost per call \\
\midrule
\textsc{EdgePotential}($u,v$) & $H \cdot n \cdot \bar{d}$ & $\mathcal{O}(1)$ \\
Mass update + bucketize       & $H \cdot n \cdot \bar{d} \cdot B$ & $\mathcal{O}(1)$ \\
Relaxation $f[h][v][b]$ update & $H \cdot n \cdot \bar{d} \cdot B$ & $\mathcal{O}(1)$ \\
Path reconstruction           & $H$ & $\mathcal{O}(1)$ \\
\bottomrule
\end{tabular}
\end{table}

\paragraph{Total time complexity.}
Summing the dominant terms in Table~\ref{tab:complexity}, the
per-decision time complexity is
\begin{equation}
T(n, e, B, H) = \mathcal{O}\!\left(H \cdot n \cdot \bar{d} \cdot B\right)
              = \mathcal{O}\!\left(H \cdot e \cdot B\right),
\label{eq:time-complexity}
\end{equation}
since $n \cdot \bar{d} = 2e$. Substituting
$H \le \alpha_{\text{budget}} \cdot \mathrm{diam}(G) \le
\alpha_{\text{budget}} \cdot n$ in the worst case yields the
loose upper bound $\mathcal{O}(n \cdot e \cdot B)$. In practice
$H$ scales with the diameter of $G$, which for the topologies in
our evaluation satisfies $\mathrm{diam}(G) = \mathcal{O}(\log n)$
for Erd\H{o}s--R\'enyi and Barab\'asi--Albert graphs and
$\mathrm{diam}(G) = \mathcal{O}(\sqrt{n})$ for the regular
lattice; the realized cost is correspondingly lower.

\paragraph{Space complexity.}
The DP stores $f$ and $\pi$ tables of dimensions
$(H+1) \times n \times B$, giving
\begin{equation}
S(n, B, H) = \mathcal{O}(H \cdot n \cdot B).
\label{eq:space-complexity}
\end{equation}
For $n = 40$ (GÉANT), $B = 32$, $H = 18$, this is approximately
$2.3 \cdot 10^4$ floating-point entries per table, or roughly
$370\,$KB across both tables in double precision. The memory
footprint is negligible for a controller-class deployment.

\paragraph{Concrete operation counts on GÉANT.}
For the headline scenario ($n = 40$, $e = 61$, $\bar{d} = 3.05$,
$H \approx 18$ for typical s--t pairs, $B = 32$):
\[
H \cdot e \cdot B \;\approx\; 18 \cdot 61 \cdot 32
                          \;=\; 35{,}136 \;\text{relaxations per
route}.
\]
With each relaxation costing on the order of 100\,ns on a
commodity CPU, this puts the per-route decision time at
approximately $3.5$\,ms, consistent with what we measure in
practice ($\sim$$3$--$8$\,ms median per route).

\paragraph{Comparison against baselines.}
The Dijkstra-based baselines (SP, LASP, LASP-aug) compute a
shortest-path tree at cost
$\mathcal{O}((n + e)\log n) \approx 700$ operations on GÉANT, or
roughly $50\times$ cheaper than EMMET-DP per decision. The
$50\times$ overhead is the price of carrying per-packet state
through a constrained dynamic program rather than a stateless
shortest-path computation. For controller-plane routing at
admission time (the deployment regime described in
the controller-plane deployment regime), this overhead is
absorbable: a $3$--$8$\,ms decision per flow is below the
inter-arrival time of new flows in any realistic backbone.
For data-plane routing at line rate, the overhead is
prohibitive without hardware acceleration, motivating the
acceleration directions discussed in the discussion.

\paragraph{Sensitivity to $B$.}
The bucket count $B$ enters linearly in
Equations~\ref{eq:time-complexity}--\ref{eq:space-complexity}.
Reducing $B$ from $32$ to $8$ (the legacy value retained as the
module default) reduces both time and memory by a factor of $4$,
at the cost of coarser mass discretization. Our hostile-audit
suspicion \#6 (\texttt{docs/AUDIT\_LOG.md})
verified that the headline GÉANT loss reduction is stable across
$B \in \{8, 16, 32, 64\}$ to within $0.20$ losses, indicating
that $B = 32$ is comfortable but not knife-edge. Practitioners
deploying EMMET-DP under tighter compute budgets can safely use
$B = 16$ with negligible quality loss.
